\documentclass[11pt]{article}
\usepackage[a4paper, top=18mm, bottom=18mm, left=20mm, right=20mm]{geometry}
\usepackage[T2A]{fontenc}
\usepackage[utf8]{inputenc}
\usepackage[ukrainian]{babel}
\usepackage{graphicx}
\setlength{\parindent}{0pt}
\setlength{\parskip}{4pt}
\pagestyle{empty}
\begin{document}
%begin
{\large\textbf{T}}

\textbf{Варіант \No\,1}\hfill \textit{ПІБ:}\,\underline{\hspace{60mm}}
\vspace{4mm}

%beginitem
1. %goodpath:Scripts/2019/Mod2019_1/Lex/01-good.txt
Відмітити все, що є однією правильною лексемою:
( 1 ) ++

( 2 ) 'c'

( 3 ) """1"2"""

( 4 ) -1.5

%enditem
%beginitem
2. %goodpath:Scripts/2018/Mod2018_1/01-good.txt
Відмітити все, що є однією правильною лексемою:
( 1 ) **

( 2 ) /

( 3 ) <=

( 4 ) 3a

%enditem
%beginitem
3. %goodpath:Scripts/2019/Mod2019_1/Lex/03-good.txt
Відмітити все, що є коректним оператором С++:
( 1 ) x=*2

( 2 ) x,y,z=z,z,y

( 3 ) r=('N'>'5')

( 4 ) 'a'a'

%enditem
%beginitem
4. %goodpath:Scripts/2018/Mod2018_1/03-good.txt
Відмітити все, що є коректним оператором С++:
( 1 ) double n.m;

( 2 ) x=int(4);

( 3 ) char c=13;

( 4 ) cout<<(3<>5);

%enditem




%end
\newpage
\end{document}


